% Contexto e instrucciones literales de las páginas 1 y 2 de EXA-0019.
% Fuente: Selectividad/matematicas_aplicadas/reserva a/Examen-Matemáticas Aplicadas a las Ciencias Sociales.pdf
% SHA-256: 6fed0384ff6fcbaaa7098ab3f35ea5f881d1280502f050ada3a7f13af528392e
% No se reproduce el logotipo ni el marco de la cabecera.
\section*{Contexto del examen original}
\textbf{PRUEBA DE ACCESO Y ADMISIÓN A LA UNIVERSIDAD}\par
CURSO 2017-2018\par
\textbf{MATEMÁTICAS APLICADAS A LAS CIENCIAS SOCIALES II}

\textbf{Instrucciones:}
\begin{enumerate}[label=\textbf{\alph*)}]
  \item Duración: 1 hora y 30 minutos
  \item Elija una de las dos opciones propuestas y conteste los ejercicios de la opción elegida.
  \item En cada ejercicio, parte o apartado se indica la puntuación máxima que le corresponde.
  \item Se permitirá el uso de calculadoras que no sean programables, gráficas ni con capacidad para almacenar o transmitir datos.
  \item Si obtiene resultados directamente con la calculadora, explique con detalle los pasos necesarios para su obtención sin el uso de la misma. Justifique las respuestas.
\end{enumerate}
\noindent\textbf{Nota de la colección:} Se reúnen las dos opciones para consulta.
La regla anterior corresponde al examen original. El curso 2017-2018 está impreso; reserva A procede sólo de la carpeta. Convocatoria y modelo no indicados. La letra de reserva no equivale a la opción A. Tabla normal auxiliar genérica asociada por carpeta.
\par\medskip
